\subsection{Mô hình dữ liệu và hợp đồng API}
\label{sec:data-api}

\subsubsection{ERD cục bộ (physical) --- Room schema v5}
\label{sec:erd-room}

Dữ liệu cục bộ của OneSpace nằm trong sáu bảng Room (schema v5). Bốn bảng tham chiếu nội dung qua mã (\texttt{videoId}, \texttt{contentId}) tới danh mục do repository cấp; đây là khóa ngoại \emph{logic}, không có ràng buộc FK vật lý vì danh mục không nằm trong CSDL quan hệ. Bảng~\ref{tab:erd-room} liệt kê đầy đủ trường và nhóm nghiệp vụ.

\begingroup
\small
\renewcommand{\arraystretch}{1.15}
\setlength{\tabcolsep}{3pt}
\begin{longtable}{>{\RaggedRight\arraybackslash}p{0.25\linewidth}>{\RaggedRight\arraybackslash}p{0.13\linewidth}>{\RaggedRight\arraybackslash}p{0.39\linewidth}>{\RaggedRight\arraybackslash}p{0.15\linewidth}}
\caption{ERD physical --- sáu bảng Room (schema v5), khóa chính gạch chân}
\label{tab:erd-room}\\
\toprule
\textbf{Bảng} & \textbf{Khóa chính} & \textbf{Trường chính} & \textbf{Nhóm nghiệp vụ} \\
\midrule
\endfirsthead
\toprule
\textbf{Bảng} & \textbf{Khóa chính} & \textbf{Trường chính} & \textbf{Nhóm nghiệp vụ} \\
\midrule
\endhead
\bottomrule
\endfoot
\texttt{video\_progress} & \underline{videoId} (FK logic) & title, positionMs (vị trí xem), updatedAt & Tương tác nội dung \\
\texttt{watchlist} & \underline{videoId} (FK logic) & title, savedAt & Tương tác nội dung \\
\texttt{downloaded\_content} & \underline{contentId} (FK logic) & type (video/nhạc/tin), title, subtitle, downloadedAt & Thư viện offline \\
\texttt{extracted\_audio} & \underline{videoId} (FK logic) & title, creator, url, duration, extractedAt & Thư viện offline \\
\texttt{point\_entries} & \underline{id} (tự tăng) & source, amount, timestamp & Gamification \\
\texttt{daily\_missions} & \underline{id} & title, targetType, target, progress, completed, date & Gamification \\
\end{longtable}
\endgroup

Ngoài Room, các thiết lập nhỏ và dữ liệu phiên nằm trong DataStore: phiên đăng nhập, cài đặt giao diện/âm thanh, bài tin đã lưu, playlists (JSON) và trạng thái game (streak, pet, shop). Tệp riêng của ứng dụng giữ các tệp âm thanh trích xuất và TTS đã tạo.

\subsubsection{Mô hình dữ liệu backend (thiết kế giai đoạn remote)}
\label{sec:erd-backend}

Bảng~\ref{tab:erd-backend} mô tả mô hình dữ liệu phía máy chủ cho giai đoạn đồng bộ theo quyết định ADR-17: backend chưa được hiện thực, danh mục hiện do repository trong app cấp. Mô hình giữ đúng các thực thể cần cho endpoint ở mục~\ref{sec:api-contract}: người dùng, nội dung, danh sách phát, điểm, nhiệm vụ và tiến độ. Ba quyết định thiết kế phục vụ đồng bộ nhiều thiết bị: (i) các bản ghi người dùng sửa được (playlist, tiến độ, nhiệm vụ) mang cột \texttt{version} để server phát hiện ghi đè và trả \texttt{409}; (ii) \texttt{point\_entries} là sổ cái chỉ ghi thêm, mỗi sự kiện có \texttt{event\_id} duy nhất nên gửi lại khi mất mạng không cộng điểm hai lần, và số dư do server tính; (iii) \texttt{sessions} chỉ lưu băm của refresh token để thu hồi được khi đăng xuất.

\begingroup
\small
\renewcommand{\arraystretch}{1.15}
\setlength{\tabcolsep}{3pt}
\begin{longtable}{>{\RaggedRight\arraybackslash}p{0.25\linewidth}>{\RaggedRight\arraybackslash}p{0.17\linewidth}>{\RaggedRight\arraybackslash}p{0.29\linewidth}>{\RaggedRight\arraybackslash}p{0.21\linewidth}}
\caption{Mô hình dữ liệu backend (thiết kế, chưa hiện thực); cột in đậm phục vụ đồng bộ}
\label{tab:erd-backend}\\
\toprule
\textbf{Entity} & \textbf{Khóa chính} & \textbf{Trường chính} & \textbf{Quan hệ} \\
\midrule
\endfirsthead
\toprule
\textbf{Entity} & \textbf{Khóa chính} & \textbf{Trường chính} & \textbf{Quan hệ} \\
\midrule
\endhead
\bottomrule
\endfoot
users & id & phone, name, avatar, created\_at & 1--n playlists, point\_entries, daily\_missions, progress, sessions \\
sessions & id & user\_id, refresh\_token\_hash, expires\_at & thuộc users \\
content\_items & id & type (video/nhạc/tin), title, url, duration, published\_at, \textbf{version} & 1--n playlist\_items, progress, quiz\_questions \\
quiz\_questions & id & item\_id, question, options, answer\_index & thuộc content\_items \\
playlists & id & user\_id, name, \textbf{version}, updated\_at & thuộc users; 1--n playlist\_items \\
playlist\_items & (playlist\_id, item\_id) & position & thuộc playlists; trỏ content\_items \\
point\_entries & id & user\_id, \textbf{event\_id} (unique), source, amount, created\_at & thuộc users (ledger điểm, server tính) \\
daily\_missions & id & user\_id, title, target, progress, date, \textbf{version} & thuộc users \\
progress & (user\_id, item\_id) & position\_ms, \textbf{version}, updated\_at & thuộc users + content\_items \\
\end{longtable}
\endgroup

\subsubsection{Hợp đồng API (thiết kế giai đoạn remote)}
\label{sec:api-contract}

Hợp đồng API dùng REST trên HTTPS với JSON và token định danh (SEC-03). Bảng~\ref{tab:api-endpoints} nhóm endpoint theo tài nguyên: Catalog phục vụ ba loại nội dung, Auth phục vụ đăng nhập OTP, Me phục vụ dữ liệu cá nhân. App hiện dùng mock repository cùng giao diện, nên thay bằng client thật không đổi tầng UI.

\begingroup
\small
\renewcommand{\arraystretch}{1.1}
\setlength{\tabcolsep}{3pt}
\begin{longtable}{>{\RaggedRight\arraybackslash}p{0.26\linewidth}>{\RaggedRight\arraybackslash}p{0.26\linewidth}>{\RaggedRight\arraybackslash}p{0.33\linewidth}>{\RaggedRight\arraybackslash}p{0.07\linewidth}}
\caption{Hợp đồng API --- endpoint chính, dữ liệu trả về và yêu cầu xác thực}
\label{tab:api-endpoints}\\
\toprule
\textbf{Endpoint} & \textbf{Mục đích} & \textbf{Trả về chính} & \textbf{Auth} \\
\midrule
\endfirsthead
\toprule
\textbf{Endpoint} & \textbf{Mục đích} & \textbf{Trả về chính} & \textbf{Auth} \\
\midrule
\endhead
\bottomrule
\endfoot
\multicolumn{4}{l}{\itshape Nhóm Nội dung (công khai, phân trang \texttt{?page=\&limit=})} \\
GET /catalog/videos & Danh mục video & id, title, creator, duration, url, thumbnail & Không \\
GET /catalog/tracks & Danh mục nhạc & id, title, artist, album, duration, url & Không \\
GET /news/articles & Danh sách tin & id, title, summary, source, publishedAt & Không \\
GET /news/articles/\{id\} & Chi tiết bài & toàn văn, quiz, bài liên quan & Không \\
\multicolumn{4}{l}{\itshape Nhóm Xác thực} \\
POST /auth/otp & Gửi mã OTP qua SMS & thời hạn 5 phút (SEC-01) & Không \\
POST /auth/verify & Đổi OTP lấy token & access token, refresh token, hồ sơ & Không \\
POST /auth/refresh & Cấp lại access token & access token mới & Refresh \\
POST /auth/logout & Thu hồi refresh token & 204 No Content & Có \\
\multicolumn{4}{l}{\itshape Nhóm Dữ liệu cá nhân (cần \texttt{Authorization: Bearer})} \\
GET /me/playlists & Playlist của tôi & id, name, items, version & Có \\
POST /me/playlists & Tạo playlist & playlist đã lưu & Có \\
PUT /me/playlists/\{id\} & Sửa playlist (kèm version) & playlist, version mới & Có \\
DELETE /me/playlists/\{id\} & Xoá playlist & 204 No Content & Có \\
GET /me/progress?since= & Tiến độ thay đổi từ mốc & danh sách itemId, positionMs, version & Có \\
PUT /me/progress/\{itemId\} & Lưu vị trí xem/nghe & bản ghi tiến độ & Có \\
GET /me/missions/today & Nhiệm vụ hôm nay & title, target, progress & Có \\
POST /me/events & Sự kiện: xong ván, xong nhiệm vụ & điểm được cộng (theo eventId) & Có \\
GET /me/points & Số dư và lịch sử điểm & balance, entries & Có \\
\end{longtable}
\endgroup

Phản hồi danh mục dùng cấu trúc phân trang thống nhất; đoạn dưới là mẫu cho \texttt{GET /catalog/videos?page=1}:

\begin{verbatim}
{ "items": [
    { "id": "v01", "title": "Giới thiệu OneSpace",
      "creator": "Nhóm OneSpace", "duration": 213,
      "url": "https://.../oceans.mp4",
      "thumbnail": "https://.../v01.jpg" } ],
  "page": 1, "hasMore": true }
\end{verbatim}

Với tài nguyên người dùng sửa được, client gửi kèm \texttt{version} đang giữ để server phát hiện ghi đè. Ví dụ sửa playlist:

\begin{verbatim}
PUT /me/playlists/p07
Authorization: Bearer <access_token>
{ "name": "Nhạc học bài", "items": ["t03", "t11"], "version": 4 }

-> 200 OK
{ "id": "p07", "name": "Nhạc học bài",
  "items": ["t03", "t11"], "version": 5,
  "updatedAt": "2026-11-07T08:00:00Z" }

-> 409 Conflict (mọi lỗi dùng cùng một dạng)
{ "error": { "code": "VERSION_CONFLICT",
    "message": "Playlist đã sửa ở máy khác",
    "details": { "serverVersion": 6 } } }
\end{verbatim}

Mọi phản hồi lỗi dùng chung cấu trúc \texttt{error\{code, message, details\}}: \texttt{code} là mã máy đọc để app chọn cách xử lý, \texttt{message} là câu hiển thị được cho người dùng, \texttt{details} mang dữ liệu bổ sung (ở đây là version hiện tại trên server để app tải lại rồi áp thay đổi). Bảng~\ref{tab:api-errors} quy định mã lỗi và cách app xử lý tương ứng. Các lỗi tầng phát và lỗi request hiện đã đi qua bộ chuyển trạng thái dùng chung (\texttt{ErrorState} kèm nút ``Thử lại''); connectivity chỉ là tín hiệu gợi ý, kết quả thao tác do lỗi thật quyết định.

\begingroup
\small
\renewcommand{\arraystretch}{1.1}
\setlength{\tabcolsep}{3pt}
\begin{longtable}{>{\RaggedRight\arraybackslash}p{0.08\linewidth}>{\RaggedRight\arraybackslash}p{0.34\linewidth}>{\RaggedRight\arraybackslash}p{0.44\linewidth}}
\caption{Mã lỗi HTTP và cách ứng dụng xử lý}
\label{tab:api-errors}\\
\toprule
\textbf{Mã} & \textbf{Ý nghĩa} & \textbf{App xử lý} \\
\midrule
\endfirsthead
\toprule
\textbf{Mã} & \textbf{Ý nghĩa} & \textbf{App xử lý} \\
\midrule
\endhead
\bottomrule
\endfoot
400 & Sai tham số & Thông báo, giữ dữ liệu cũ \\
401 & Token hết hạn & Yêu cầu đăng nhập lại \\
403 & Không có quyền & Ẩn hành động cần quyền \\
404 & Không có nội dung & EmptyState + gợi ý thay thế \\
409 & Xung đột phiên bản & Giữ bản mới hơn, đồng bộ lại \\
429 & Quá nhiều request & Chờ rồi tự thử lại \\
5xx & Lỗi server & ``Không tải được'' + Thử lại \\
\end{longtable}
\endgroup
